
\section{Structural properties}\label{pts:metatheory}

The following facts let us use finite open contexts and chosen typing
derivations without assuming uniqueness of sorts or product rules.
These are standard properties of pure type systems; compare
\cite[Chapter~2]{Poll1994}. We give the proofs needed here, especially the
steps concerning annotations and nonfunctional specifications.

\begin{lemma}[Validity and scope]\label{pts:validity}
If $\Gamma\vdash M:A$, then $\Gamma$ is valid, and
$\mathrm{FV}(M)\cup\mathrm{FV}(A)\subseteq\operatorname{dom}(\Gamma)$.
For each declaration of $\Gamma$, its sorting judgment in the preceding
prefix occurs as a subderivation of the given derivation.
\end{lemma}

\begin{proof}
Induct on the derivation. An axiom has empty context. The variable and
weakening rules supply the sorting judgment for their last declaration;
the induction hypothesis supplies the earlier ones. Every other rule
has a premise with unchanged context, so its induction hypothesis
supplies validity and the required subderivations. Scope follows in the
same induction. For application, use
\[
 \mathrm{FV}(B[x:=N])\subseteq
       (\mathrm{FV}(B)\setminus\{x\})\cup\mathrm{FV}(N).
\]
The binder cases remove their fresh bound variable from the free
variables of the body.
\end{proof}

\begin{lemma}[Thinning by insertion]\label{pts:thinning}
Suppose $\Gamma\vdash M:A$, and suppose that $\Delta$ is a valid context
containing the declarations of $\Gamma$ unchanged and in their original
order. Then $\Delta\vdash M:A$.
\end{lemma}

\begin{proof}
Induct on the size of the given derivation. For an axiom, weaken through
the declarations of $\Delta$. For a variable or weakening rule, split
$\Delta$ just after the last declaration of the source context.
Transport the rule's premises to the part before that declaration by
induction, reapply the rule, and then weaken through the remaining
suffix. The necessary sorting judgments for that suffix are supplied
by validity of $\Delta$.

At a binder, first rename its bound variable fresh for $\Delta$.
Its domain sorting is a premise of product formation or a smaller
subderivation supplied by Lemma~\ref{pts:validity}. Induction transports
this sorting to $\Delta$, making the extended target context valid.
Induction therefore also transports the body or codomain premise to
that extension. Reapply the original rule. Application and conversion
transport their unchanged-context premises and reapply their rule.
\end{proof}

\begin{lemma}[Typed substitution]\label{pts:substitution}
Let $\Gamma=(x_1:A_1,\ldots,x_n:A_n)$ and let $\Delta$ be valid.
Suppose a simultaneous substitution $\sigma$ satisfies
\[
 \Delta\vdash\sigma(x_i):A_i[\sigma|_{\{x_1,\ldots,x_{i-1}\}}]
 \qquad(1\leq i\leq n).
\]
Then $\Gamma\vdash M:A$ implies $\Delta\vdash M\sigma:A\sigma$.
In particular, if
\[
 \Gamma,x:D,\Theta\vdash M:A,
 \qquad \Gamma\vdash N:D,
\]
then
\[
 \Gamma,\Theta[x:=N]\vdash M[x:=N]:A[x:=N].
\]
\end{lemma}

\begin{proof}
Raw substitution preserves beta conversion. Indeed, a substituted root
contraction is a contraction after renaming bound variables fresh;
the required equality is the composition law for capture-avoiding
substitutions. Compatibility extends this to every reduction position,
and hence to conversion chains.

For the simultaneous assertion, induct on the source derivation.
Thin axioms to $\Delta$. At a variable use its assigned judgment, and
at weakening discard the assignment for the last declaration. Under a
binder $x:D$, choose $y$ fresh for $\Delta$ and the images of $\sigma$.
Transport the domain sorting by induction, extend the target context
to $\Delta,y:D\sigma$, and extend the substitution by $x\mapsto y$.
The earlier assignments thin to this context; the new assignment is
the variable rule. Apply induction to the binder premise and reapply
product formation or abstraction with the original axiom or rule
choices. Application uses substitution composition. Conversion uses
the transported sorting of its target and preservation of raw
conversion.

For single substitution with a dependent suffix, construct the target
context one declaration at a time. Start with the identity assignments
on $\Gamma$ and the image $x\mapsto N$. Transport the first suffix
declaration's sort judgment by the simultaneous assertion, append that
declaration, thin the old assignments, and assign its new variable to
itself. Repeat through $\Theta$, then transport the desired judgment.
This also proves validity of the substituted target context.
\end{proof}

\begin{lemma}[Conversion of a declaration]\label{pts:context-conversion}
Suppose $\Gamma\vdash D:s$, $\Gamma\vdash D':s'$, and $D=_\beta D'$.
Replacing $x:D$ by $x:D'$ in any valid context
$\Gamma,x:D,\Theta$ preserves validity and every derivable judgment,
with the subject, expected type, and later declarations unchanged.
The converse replacement has the same property.
\end{lemma}

\begin{proof}
In $\Gamma,x:D'$ the variable rule gives $x:D'$. Thin the old sorting
of $D$ to this context and convert to obtain $x:D$. Thus the identity
on the old variables is a typed substitution from $\Gamma,x:D$ to
$\Gamma,x:D'$. As in the suffix construction in
Lemma~\ref{pts:substitution}, transport each later sorting judgment,
append the same declaration, and extend the identity substitution.
Finally transport the desired judgment. Interchanging $D,D'$ proves
the reverse assertion. Nothing requires $s=s'$.
\end{proof}

\Needspace{4\baselineskip}
\begin{lemma}[Generation]\label{pts:generation}
For derivable judgments in an arbitrary pure type system:
\begin{enumerate}
\item If $\Gamma\vdash x:T$, then $x:D$ occurs in $\Gamma$ and
      $T=_\beta D$.
\item If $\Gamma\vdash s:T$, then $(s,u)\in\mathcal A$ and
      $T=_\beta u$ for some sort $u$.
\item If $\Gamma\vdash\Pi x:D.E:T$, then for some
      $(a,b,c)\in\mathcal R$,
      \[
       \Gamma\vdash D:a,\qquad
       \Gamma,x:D\vdash E:b,\qquad T=_\beta c.
      \]
\item If $\Gamma\vdash\lambda x:D.M:T$, then for some $E,c$,
      \[
       \Gamma,x:D\vdash M:E,\qquad
       \Gamma\vdash\Pi x:D.E:c,\qquad
       T=_\beta\Pi x:D.E.
      \]
\item If $\Gamma\vdash MN:T$, then for some $D,E,x$,
      \[
       \Gamma\vdash M:\Pi x:D.E,\qquad
       \Gamma\vdash N:D,\qquad T=_\beta E[x:=N].
      \]
\end{enumerate}
Every immediate syntactic child of a typable expression is typable
in the corresponding context, with a binder declaration added for a
body or codomain.
\end{lemma}

\begin{proof}
Induct on the derivation, stripping any final conversion and weakening
steps. Conversion composes the asserted conversion of expected types.
For weakening, thin all extracted premises back to the original
context, including under a fresh binder. When neither rule is last,
the outer syntax identifies its introducing rule and supplies the
displayed premises. The child assertion follows from these premises;
for a lambda's annotation use validity of the body context.
\end{proof}

\begin{lemma}[Correctness of types]\label{pts:correctness}
If $\Gamma\vdash M:A$, then either $A$ is a literal sort or
$\Gamma\vdash A:s$ for some sort $s$. Moreover, every type assigned
to an application is sorted, including an assigned type that is itself
a literal sort. Consequently every nonliteral legal expression is
typable.
\end{lemma}

\begin{proof}
Induct on the derivation. The axiom and product conclusions have literal
sorts as expected types. Variable, weakening, abstraction, and conversion
use their sorting premises, thinning where necessary. In the application
case, induction on the function premise sorts its product type, since
that type is not a literal sort. Product generation then gives a sorting
of the codomain under the binder. Substitution sorts the instantiated
result type. This also proves the stronger application assertion at
a direct application rule; final weakening preserves it by thinning,
and final conversion explicitly sorts the new expected type. Every
legal expression that occurs as an expected type is therefore either
typable or a literal sort.
\end{proof}

\subsection{Conversion and preservation of typing}
\label{pts:reduction-metatheory}

The next lemma concerns raw expressions. Thus it applies before any
normalization assumption, regardless of which products a specification
can form.

\begin{lemma}[Confluence and product compatibility]\label{pts:confluence}
Full beta reduction on raw expressions is confluent. Therefore distinct
sorts are not convertible, no product is convertible to a sort, and
\[
 \Pi x:D.E=_\beta\Pi x:D'.E'
 \quad\Longrightarrow\quad
 D=_\beta D'\ \hbox{ and }\ E=_\beta E',
\]
after consistently renaming the bound variables.
\end{lemma}

\begin{proof}
Define parallel contraction $\Rightarrow$ by reflexive variable and
sort clauses, homomorphic clauses for all expression constructors, and
the additional clause
\[
 \frac{D\Rightarrow D'\qquad M\Rightarrow M'\qquad N\Rightarrow N'}
      {(\lambda x:D.M)N\Rightarrow M'[x:=N']}.
\]
In particular the homomorphic binder clauses reduce both children.
Induction on a parallel derivation proves parallel substitution:
if $M\Rightarrow M'$ and $N\Rightarrow N'$, then
$M[x:=N]\Rightarrow M'[x:=N']$. For its contracting case, rename
the contracted binder fresh and use composition of substitutions;
the other cases reapply the corresponding homomorphic clause.

Define $M^*$, its complete development, recursively by developing every
child, with
\[
 ((\lambda x:D.P)Q)^*=P^*[x:=Q^*]
\]
at an application whose original function is a lambda. At any other
application use $(PQ)^*=P^*Q^*$. We claim that
$M\Rightarrow N$ implies $N\Rightarrow M^*$. Prove this by induction
on the parallel derivation. At an original root redex, either the
first parallel step was homomorphic, in which case contract the root
in the second step, or it contracted the root, in which case parallel
substitution applies. At an original nonredex application, use the
homomorphic clause in the second step, even if its function has become
a lambda. The remaining cases follow by the corresponding
homomorphic clauses and induction. Thus any two parallel reducts of
$M$ parallel-reduce to $M^*$, proving the diamond property.

Every ordinary step is parallel. Conversely, every parallel step is
a finite sequence of ordinary steps: first reduce the children and
then perform its root contraction, if any. Hence the reflexive-transitive
closures of the two relations coincide, and beta reduction is confluent.
Convertible expressions consequently have a common reduct. Reducing
a product never removes its outer constructor, whereas a sort is
irreducible. Comparing the common reducts gives all the asserted
compatibility and separation properties.
\end{proof}

\begin{lemma}[Subject reduction]\label{pts:subject-reduction}
If $\Gamma\vdash M:A$ and $M\to_\beta^*M'$, then
$\Gamma\vdash M':A$, with the exact original expected type.
\end{lemma}

\begin{proof}
It suffices to consider one step. Induct on the typing derivation.
There are no steps from sorts or variables. At a final weakening or
conversion, apply induction to its subject premise and reapply the
same rule.

At product formation, a step in the domain preserves its original sort
by induction. Lemma~\ref{pts:context-conversion} transports the codomain
sorting to the changed binder context, and the same product rule
re-forms the result. A codomain step is handled by induction on its
typing premise.

At abstraction, a step in the body is immediate by induction. For an
annotation step $D\to_\beta D'$, the supporting judgment
$\Gamma\vdash\Pi x:D.E:c$ is a strictly smaller premise derivation.
Apply induction to its corresponding domain step to get
$\Gamma\vdash\Pi x:D'.E:c$. Product generation sorts $D'$;
context conversion then transports the body judgment to
$\Gamma,x:D'$. Abstraction gives the changed lambda its new product
type. Convert back to the old product type, which is sorted by the
original supporting premise.

At application, a step in the function preserves its product type by
induction. A step $N\to_\beta N'$ in the argument gives the new
result type $E[x:=N']$; convert this to $E[x:=N]$, whose sorting was
proved in Lemma~\ref{pts:correctness}.

It remains to treat a root contraction
$(\lambda x:D_0.P)N\to_\beta P[x:=N]$. Suppose the application
premises assign the function type $\Pi x:D.E$ and the argument
type $D$. Generation of the function judgment gives some $E_0,c$
such that
\[
 \Gamma,x:D_0\vdash P:E_0,
 \qquad\Gamma\vdash\Pi x:D_0.E_0:c,
 \qquad\Pi x:D_0.E_0=_\beta\Pi x:D.E.
\]
Product compatibility gives $D_0=_\beta D$ and $E_0=_\beta E$.
Generation of the supporting product sorts $D_0$, so convert
$N:D$ to $N:D_0$. Substitution yields
$\Gamma\vdash P[x:=N]:E_0[x:=N]$.
Convert to the old application type $E[x:=N]$, which is sorted by
Lemma~\ref{pts:correctness}. This completes every reduction position,
including positions inside annotations.
\end{proof}

\subsection{Neutral expressions and normal forms}
\label{pts:neutrals}

An expression $xN_1\cdots N_k$, with $k\geq0$, is called
\emph{neutral}. Its arguments need not be normal. This restricted
class has unique types even when the specification does not.

\begin{lemma}[Uniqueness for neutral expressions]\label{pts:neutral-uniqueness}
If $H$ is neutral and $\Gamma\vdash H:A$ and $\Gamma\vdash H:B$,
then $A=_\beta B$. If $\Gamma\vdash H:s$ for a literal sort $s$,
then $s$ itself is sorted in $\Gamma$. Every normal typable application
is neutral.
\end{lemma}

\begin{proof}
Induct on the length of the neutral spine. At a variable, generation
identifies both types with its declaration type up to conversion.
For $HN$, generation of its two typings gives function product types
$\Pi x:D.E$ and $\Pi x:D'.E'$. Induction makes these convertible.
Product compatibility gives $E=_\beta E'$, so substitution makes
their instantiated codomains convertible. These codomains are
convertible to the original two result types by generation.

If a variable has type $s$, its declaration type $D$ is sorted by
validity and thinning, and $D=_\beta s$ by generation. Confluence
and irreducibility of $s$ give $D\to_\beta^*s$. Subject reduction
therefore sorts $s$. For a nonempty neutral spine, the stronger
application assertion in Lemma~\ref{pts:correctness} sorts its
expected type directly.

Finally, the head of a normal application cannot be a lambda, since
that would give a redex. It cannot be a sort or product: generation
makes any type of either head convertible to a sort, whereas the
first application requires a product type. Product/sort separation
excludes this. The only remaining head is a variable.
\end{proof}

\begin{remark}\label{pts:nonfunctionality}
General uniqueness of types is false under our hypotheses. For example,
if $(r,a)$ and $(r,b)$ are axioms with distinct $a,b$, then $r$ has
the two nonconvertible types $a,b$. Lemma~\ref{pts:neutral-uniqueness}
does not apply to the sort $r$. No later use of uniqueness concerns
an arbitrary expression.
\end{remark}

\begin{lemma}[Normalizing judgments]\label{pts:normalizing-judgments}
Assume that every legal expression is weakly normalizing. Write
$M^\#$ for its normal form, uniquely determined up to alpha equivalence.
If $\Gamma\vdash M:A$, then
\[
                \Gamma\vdash M^\#:A
                \qquad\text{and}\qquad
                \Gamma\vdash M^\#:A^\#.
\]
The second judgment is also available with $M$ in place of $M^\#$.
If $\Gamma\vdash A:s$, then $\Gamma\vdash A^\#:s$.
A sorted normal expression is either a sort, a neutral expression,
or a product.
\end{lemma}

\begin{proof}
Weak normalization supplies normal forms; confluence gives uniqueness.
Subject reduction proves the first judgment and the assertion about
sorting. Correctness of types says that $A$ is sorted or is a literal
sort. In the first case subject reduction sorts $A^\#$, allowing
conversion of the expected type. In the second case $A^\#=A$, so
there is nothing to convert if that sort is unsorted. Finally a
normal application is neutral by Lemma~\ref{pts:neutral-uniqueness},
and a lambda cannot have a sort type by generation and product/sort
separation.
\end{proof}

\subsection{Finite specifications and a common context}
\label{pts:finite-and-context}

For the main implication it is enough to work with a finite
specification. We will then place all finite contexts inside one
countable context, while each judgment continues to use only a finite
prefix.

\begin{lemma}[Finite restriction]\label{pts:finite-restriction}
Let $\mathcal P=(\mathcal S,\mathcal A,\mathcal R)$ be a pure type
system.  An expression is called \emph{legal} if it occurs as the subject
or the expected type of a derivable judgment.  If every legal expression
of $\mathcal P$ is weakly normalizing for full $\beta$-reduction, but some
legal expression is not strongly normalizing, then there are finite sets
\[
  \mathcal S_0\subseteq\mathcal S,\qquad
  \mathcal A_0\subseteq\mathcal A\cap\mathcal S_0^2,\qquad
  \mathcal R_0\subseteq\mathcal R\cap\mathcal S_0^3
\]
such that the restricted system
$\mathcal P_0=(\mathcal S_0,\mathcal A_0,\mathcal R_0)$ has the same two
properties.
\end{lemma}

\begin{proof}
First recall how legality relates to typability.  If $E$ is a subject,
it is typable by definition.  If $E$ is an expected type, correctness of
types says that either $E$ is a literal sort or $\Gamma\vdash E:s$ for
some context $\Gamma$ and sort $s$.  Thus every nonliteral legal
expression is typable.  Also, context validity says that each declaration
type is sorted in its preceding context, so declaration types introduce
no further kind of untypable expression.  Literal sorts have no
$\beta$-redexes.

Choose a legal expression $E$ that is not strongly normalizing.  It is
not a literal sort, so fix a finite derivation of a judgment
$\Gamma\vdash E:B$.  For every use of conversion in this derivation,
choose a finite raw $\beta$-conversion chain witnessing its side
condition.  Let $\mathcal S_0$ contain all sorts occurring in the
derivation and these chains.  Let $\mathcal A_0$ and $\mathcal R_0$
contain the axiom pairs and formation triples used in the derivation,
enlarging $\mathcal S_0$ to contain their entries if necessary.  These
sets are finite, and the same derivation and conversion chains establish
$\Gamma\vdash E:B$ in $\mathcal P_0$.

A full $\beta$-step introduces no new sort symbol: substitution can only
copy symbols already present in the redex, and contextual reduction has
the same property, including reduction in annotations.  Consequently
the infinite reduction from $E$ witnessing failure of strong
normalization is also a reduction in the raw syntax over
$\mathcal S_0$.

Conversely, every judgment derivable in $\mathcal P_0$ is derivable in
$\mathcal P$.  If $F$ is legal in $\mathcal P_0$, the weak-normalization
hypothesis therefore supplies a finite reduction $F\to_\beta^*N$ in
$\mathcal P$ with $N$ normal.  No term in this reduction contains a sort
outside $\mathcal S_0$.  The reduction is thus available in
$\mathcal P_0$, and normality is unchanged, since it is the absence of a
syntactic $\beta$-redex.  This proves weak normalization of every legal
expression of $\mathcal P_0$, while $E$ remains a legal expression that
is not strongly normalizing.
\end{proof}

\begin{lemma}[A saturated context]\label{pts:saturated-context}
Suppose that the sort set $\mathcal S$ is at most countable, and that
the variable set is countably infinite.  There is a possibly empty
countable sequence of declarations $\Omega$ with the following
properties.  A judgment over $\Omega$ means a judgment derivable over
some finite initial segment of $\Omega$.
\begin{enumerate}
\item Every finite initial segment of $\Omega$ is valid.  An infinite
set of variable names is disjoint from all its declared names.
\item If $P$ is a finite initial segment of $\Omega$ and
$P\vdash D:s$, then $\Omega$ contains infinitely many declarations
$y:D$ after $P$, with distinct declared names.
\item Every finite valid context embeds, by an injective renaming of
its declared variables, as a subsequence of a finite initial segment
of $\Omega$.  If $\Gamma\vdash M:A$, such an embedding $\rho$ gives
$P\vdash M\rho:A\rho$ for a finite initial segment $P$ containing the
embedded context.
\item If $P\vdash D:s$ for a finite initial segment $P$, and $F$ is any
finite set of variable names, there is a larger finite initial segment
$Q$ ending in a declaration $y:D$, where
$y\notin F\cup\operatorname{dom}(P)$, and $Q\vdash y:D$.
\end{enumerate}
Moreover, finitely many judgments over $\Omega$ hold over a common
finite initial segment.  Injective renaming in the third clause
preserves and reflects weak and strong normalization.
\end{lemma}

\begin{proof}
Partition the variable names into disjoint infinite sets $V_0$ and
$V_1$.  Only names in $V_0$ will be declared in $\Omega$; names in
$V_1$ remain available for bound variables.  Because $\mathcal S$ and
the variable set are countable, the set of finite raw expressions is
countable.  We work up to $\alpha$-equivalence, or equivalently choose
representatives when enumerating expressions.

Choose a sequence of pairs $(p,D)$, with $p$ a nonnegative integer and
$D$ a raw expression, in which every pair occurs infinitely often.
Construct finite contexts $\Delta_n$ starting with the empty context.
At stage $n$, consider the scheduled pair $(p,D)$.  If $\Delta_n$ has
at least $p$ declarations and its initial segment of length $p$ derives
$D:s$ for some sort $s$, append a declaration $y:D$, choosing
$y\in V_0\setminus\operatorname{dom}(\Delta_n)$.  Otherwise leave
$\Delta_n$ unchanged.  This is a set-theoretic construction and does
not require a decision procedure for typing.  In the append case,
scope gives $\mathrm{FV}(D)\subseteq
\operatorname{dom}(\Delta_n)$, and thinning gives
$\Delta_n\vdash D:s$.  Thus the new declaration is fresh and the
extended context is valid.  Induction proves validity at every stage.

Let $\Omega$ be the sequence obtained by retaining all declarations
appended in this construction.  Every finite initial segment appears
at some stage and is valid.  No name from $V_1$ is declared.  If
$P\vdash D:s$ and $P$ has length $p$, then, once $P$ has been
constructed, every later occurrence of the pair $(p,D)$ appends a new
declaration of type $D$.  There are infinitely many such occurrences.
This proves the first two clauses, including the case where no
declaration is ever possible and $\Omega$ is empty.

For the embedding clause, write a finite valid context as
\[
  \Gamma=(x_1:D_1,\ldots,x_m:D_m).
\]
Induct on its prefixes.  Suppose an injective renaming $\rho_i$ embeds
the first $i$ declarations, in their original order, into a finite
initial segment $P_i$ of $\Omega$.  Validity supplies
\[
  x_1:D_1,\ldots,x_i:D_i\vdash D_{i+1}:s
\]
for some sort $s$.  Renaming and then thinning yield
$P_i\vdash D_{i+1}\rho_i:s$.  By the second clause, a later
declaration has the form $y:D_{i+1}\rho_i$, with $y$ fresh for the
already chosen names.  Set $\rho_{i+1}(x_{i+1})=y$ and take an initial
segment through this declaration.  This retains all previously
embedded declarations unchanged and in order.  The induction begins
with the empty context and proves the desired embedding.

Renaming a derivation gives
$\Gamma\rho\vdash M\rho:A\rho$.  The chosen initial segment is a
valid extension by insertion of $\Gamma\rho$, so thinning gives the
asserted judgment over that initial segment.  All substitutions and
renamings here are capture-avoiding; the reserved names permit bound
variables to be renamed away from the declarations whenever needed.

For the fourth clause, only finitely many of the infinitely many
declarations of type $D$ after $P$ can have their declared names in
$F$.  Choose any other one and let $Q$ be the initial segment ending
there.  Thinning the judgment for $D$ to the prefix immediately before
that declaration, followed by the variable rule, gives $Q\vdash y:D$.
In particular, when finitely many substitution images have their free
variables in $P$, this $y$ is fresh for all those images.

Finally, take the longest among finitely many initial segments
witnessing judgments over $\Omega$ and thin all the judgments to it.
An injective renaming of the finitely many variables of a context can
be extended to a bijective renaming of the whole variable set.  Such
a renaming commutes with full $\beta$-reduction and its inverse does
too.  It therefore preserves and reflects normal forms, finite
normalizing reductions, and infinite reductions.  This proves the
last assertions.
\end{proof}
